Multi-touch gesture-based gaming is a relatively unexplored topic. This means that there is a considerable lack of existing information on the way in which interaction techniques influence the gaming experience on such devices. Given the recent surge in commercially available touch-devices, in addition to the fact that this technology may provide a more immersive gaming environment, it appears that research in this area may produce valuable information related to improving gaming experience. This project aimed to address this issue by investigating the effect of direct manipulation and gestures on the level of flow elicited during game-play. In addition, an interface using a mix of these interaction techniques was constructed in order to assess their combination. Additional elements, such as game genre, as well as previous gaming or touch-device experience and their effects on game enjoyment in relation to interaction style, were also investigated.

In order to conduct this study, a number of components were developed. These included a feature extraction module, two separate gesture recognition engines as well as three games, each of a different genre. These components were integrated for the purpose of user experimentation.

Experimentation revealed a significant difference between the interaction styles. That is, the gestural interface was found to elicit a higher degree of flow than either of the other two interaction techniques. Furthermore, as the mixed interaction style game was more successful in creating a flow state than the direct manipulation version, but only in conditions where the same was true of the gestural interface, one might conclude that this was due to the inclusion of gestural commands within the mixed interface. The direct manipulation interface was found to be more intrinsically rewarding, however. This may be due to the increased cognitive load placed on the user during game-play, as memorisation and recall of gestures is required. Overall, however, the gestural interface appears to elicit the highest degree of flow, when compared to either of the other two interfaces. This was found to be consistent across game genre and previous gaming and touch-device experience.

While game genre and specific participant characteristics do not appear to influence which interface version is preferred, both variables show variation when considering the degree to which each interface elicits flow. Variations were found between the level of flow experienced when playing games of different genres. The Strategy and Fighting/Duelling games both appear to have produced more flow than the Puzzle game. This variation, however, was relatively consistent across interface type, suggesting that some genres may, in general, elicit flow to a greater or lesser extent. Similar results were present when considering the effect of previous gaming and touch-device experience. Individuals who lacked either significant gaming or touch-device experience appear to have experienced less of a flow state than individuals who possess experience in these areas, regardless of interface type. Gaming experience appears to be a stronger predictor for this effect than touch-device experience. In addition, gaming experience seems to suggest a greater sensitivity to overall flow trends, with individuals experiencing more extreme flow states both positively and negatively. This may be due to the fact that previous gaming experience allows individuals to more easily enter a flow state. The difference between those with gaming or touch-device experience and those without may also be due to a lack of familiarity with either touch-based interaction, or the style of application in general. Overall, however, genre and previous gaming and touch-device experience should be viewed as predictors for the degree to which flow can be experienced, rather than indicators of which interface version is most enjoyable.

A number of limitations within the present study should be noted, however. The first of these is the small sample size when considering the number of variables. That is, for the genre and participant characteristic variables, only 15 participants per condition were used. This is further compounded by the fact that the measure used assessed flow along nine separate dimensions. Thus, one should be cautious when attempting to generalise the findings of this study to a larger population. In addition, as computer literacy was controlled for, the generalisability of the study is further limited. Furthermore, as only one game from each genre was tested, the results related to genre may be due to the characteristics of the individual games rather than being indicative of that genre in general. The specific nature of the games tested can be expanded on further in this regard. That is, as the games were designed with gestural interaction in mind, this may have influenced the results.

Thus, although gestures were found to elicit more flow, they should not be seen as a panacea. Rather, the decision to use gestures rather than some other interaction technique for multi-touch gaming should be consciously examined within the context of the specific game being developed. While gestures may prove ideal for some games, their use and associated cognitive load, may detract from the overall flow state achieved in other games. With this said, the results seem to indicate that for games where their use is appropriate, the inclusion of gestural controls when compared to other interaction techniques, may increase the degree of flow, and thus overall game enjoyment, experienced by players.

\section{Future work}
A number of elements within this project could be expanded on further in future work. The first of these relates to the sample size used for experimentation. A similar experiment with a larger sample size would provide more conclusive results. In addition, it would allow for the investigation of additional subject characteristics, such as computer literacy, with regard to their effect on the degree of flow experienced when using different interaction techniques. Furthermore, the addition of a pre-experiment questionnaire used to assess an individual's propensity to experience flow could also provide valuable information with regard to subject characteristics. This would also help to differentiate between gaming experience and flow propensity.

An investigation into the effects of game genre on interaction technique preference would also benefit from a larger sample size. This is due to the fact that the existing sample size limited the number of games it was possible to test within each genre. The inclusion of multiple games per genre, as well as additional genres, could provide more generalisable results in this regard. Further research could also examine the reasons for the differences in the amount of flow elicited by each interaction technique across the various flow dimensions. In particular, further study into why the direct manipulation version elicits more flow along the \emph{Autotelic Experience} dimension than the other interfaces may prove valuable. This could be expanded further to include an investigation of whether there are specific gestural or direct manipulation components that can be shown to increase the possibility of experiencing a flow state. In addition, means of reducing the cognitive load placed on users by gestures could be investigated within the context of computer games.

Overall, however, similar studies with a wider range of games and a larger sample size could greatly improve the generalisability of this work. In addition, the investigation of further variables related to this concept would help to determine more conclusively when the use of gestures is most appropriate in multi-touch gaming.
